% appendix_cards.tex: the revolving-credit extension: the design (moved from Section 6.1) and its feasibility in the archives.
\section{Revolving Credit: the Design and Its Feasibility in the Archives}\label{appendix:cards}

\subsection{The design}\label{sec:cards}

The estimator applies to revolving credit, where no contract fixes the payment schedule, once the two margins are defined by the account's position rather than by its contract. Each month a cardholder pays an amount between the issuer's formula minimum, $m = \alpha B$ for a balance $B$ and a formula rate $\alpha$, and the full balance, or defaults. Three positions are possible. A borrower who pays in full faces no binding minimum and has a default probability near zero. A borrower who pays more than the minimum but less than the balance satisfies the Euler equation at the card's interest rate; the minimum is slack and default does not respond to it at the margin. A borrower who pays exactly the minimum would pay less if permitted: the Euler equation at the card rate fails, and default is the margin on which the minimum acts. Proposition \ref{prop:sensT} applies to the third group, whose membership is observable from the payment record. Borrowers in the first two positions contribute nothing to the payment margin; the one exception, a wealth effect of the balance on interior revolvers, is removed by restricting the sample to observed minimum-payers.

For a minimum-payer the payment stream is geometric rather than level, since the formula recomputes the minimum on the declining balance, and its duration is the horizon over which the payment stream runs off. At the median formula rate observed in the archives, $0.036$ to $0.038$ of the balance, and an annual percentage rate of twenty percent the duration is $3.8$ to $4.3$ years (Table \ref{tab:cards_probe}), which places revolving credit near auto loans on the horizon axis, not at its far end. Because the horizon depends on $\alpha$ and the interest rate and not on $B$, the payment and the horizon map onto observables in a way the installment market does not offer: a change in the balance at a fixed formula moves the payment with the horizon held fixed, which is the payment margin, and a change in the formula moves the payment and the horizon together, with the horizon's response known from the annuity identity; and the expiration of a promotional rate moves the minimum payment at a date known in advance at a fixed balance, the variation of Proposition \ref{prop:curveT}. Each margin has a canonical source of variation in the literature, credit-limit discontinuities at score cutoffs for the balance \citep{agarwal_credit_2018, gross_liquidity_2002} and issuer minimum-payment rules for the formula \citep{keys_minimum_2019}, and neither relies on cross-sectional variation in the minimum that the borrower's own balance choice contaminates. Survival $Q_s$ is the probability of still paying the minimum at $s$, with exits by payoff and by default both counted, and it is measurable in the panel.

Two qualifications define what the exercise identifies. Being at the minimum is chosen, on preferences, income, and wealth, exactly as being an auto borrower is chosen; a card estimate, once the formula and rate changes the design requires are measured, would be the discount rate of borrowers at the card's default margin and comparable to the auto estimate on that basis, not a rate for cardholders at large. And a change in the formula or the limit that moves a borrower across positions mixes a change of position with a response at the margin; the first-stage effect of each instrument on position is the check. The archives used in this paper contain the fields the design needs for revolving trades, and Appendix \ref{appendix:cards} reports the feasibility probe; the promotional-rate expiration, a known change in the payment at a known future date for borrowers who are still revolving, is the version of the Proposition \ref{prop:curveT} experiment that this market uniquely offers.

\subsection{Feasibility}
The design needs, for revolving trades, the balance, the formula minimum, the payment actually made, the manner-of-payment rating, and the furnisher. The trade archives contain each of these fields. Table \ref{tab:cards_probe} reports, for two archive months, the share of revolving trades with a positive balance whose payment equals, exceeds, or falls short of the minimum; the distribution of the formula rate $\alpha = m/B$; and the duration of the minimum-payment stream at the observed $\alpha$ and a twenty percent annual rate. (\texttt{Code/extract/E5\_cards\_probe.R}.)

\begin{table}[ht]
\centering
\caption{Revolving trades in the archives: payment position, formula rate, and horizon}
\label{tab:cards_probe}
% replication: Code/extract/E5_cards_probe.R, Code/exhibits/make_placeholder_tables.R; results/cards_probe.csv
\begin{tabular}{lrr}
\toprule
 & 201807 & 202404 \\
\midrule
Revolving trades & 54,633,873 & 62,516,556 \\
Trades with a computable formula rate & 48,437,666 & 53,725,970 \\
Share paying the minimum & 0.131 & 0.159 \\
Share paying above the minimum & 0.772 & 0.769 \\
Share paying below the minimum & 0.096 & 0.072 \\
Formula rate, median & 0.0357 & 0.0380 \\
Formula rate, p10--p90 & 0.0149--0.2358 & 0.0157--0.2358 \\
Duration at the median formula rate (years) & 4.3 & 3.8 \\
\bottomrule
\end{tabular}

\end{table}
